%% hopfion-framework/profile/virial.tex
%% Sources: main_paper1.tex, main_paper2.tex, main_paper3.tex
%% Virial condition, saddle point b*, Newton series for q*, beta* derivation
%% Auto-generated from project files. Edit source papers to update.


%════════════════════════════════════════════════════════════════════
% Virial Condition and beta* — Paper I
%════════════════════════════════════════════════════════════════════

\begin{definition}[Saddle snapshot]
\label{P1:def:snapshot}
The \emph{saddle snapshot} is the configuration reached by gradient
descent on $E_{\mathrm{geom}}$ at the step where the Derrick scale factor
\begin{equation}
  \sopt \;\equiv\; \sqrt{\frac{\lam\,\Jfour}{\Ja+\ms\,\Jiso}}
  \label{P1:eq:sopt_def}
\end{equation}
is closest to $1$.  At this snapshot the Derrick balance is satisfied
and the configuration is the saddle of the full energy $E$.
\end{definition}
\status{published}
\source{P1:def:snapshot}

\begin{definition}[Self-consistent coupling]
\label{P1:def:bstar}
$\bst$ is the unique value at which the EL minimum of $\Efb$ in the
$Q=2$ sector satisfies
\begin{equation}
  \frac{\Jfb}{\Ja} \;=\; \vp.
  \label{P1:eq:virial_def}
\end{equation}
Numerically, $\bst \approx 0.452$.
\end{definition}
\status{published}
\source{P1:def:bstar}

\begin{proposition}[$P_{\Kfb} = 2\vp\,P_{\Ja}$ is an algebraic identity]
\label{P1:prop:PKfb}
Let $P_I = d I[f(r(1+\varepsilon),z)]/d\varepsilon|_0$ denote the
variation of functional $I$ under the radial dilation
$r\mapsto r(1+\varepsilon)$, $z$ and $R_0$ fixed.
At the self-consistent coupling $\bst$ satisfying
$\Kfb = \Ja + \ms\Jfb = 2\vp\Ja$ and $\ms = 3-\vp$:
\begin{equation}
  P_{\Kfb} \;=\; 2\vp\,P_{\Ja}.
  \label{P1:eq:PKfb_exact}
\end{equation}
\end{proposition}
\status{published}
\source{P1:prop:PKfb}

\begin{proposition}[Feedback stiffness]
\label{P1:prop:Kfb}
At the self-consistent $\bst$,
$\Kfb = 2\vp\,\Ja$.
\end{proposition}
\status{published}
\source{P1:prop:Kfb}

\begin{lemma}[Virial sum]
\label{P1:lem:virial}
$\vp^{-1} + \ms = 2$ exactly.
\end{lemma}
\status{published}
\source{P1:lem:virial}

\begin{conjecture}[Linking Scale Conjecture --- proved in Paper~II]
\label{P1:conj:linking}
For the $Q$-Hopfion at WZW level $k$, the EL saddle of $\Efb$
in the topological sector with virial condition
$\Jfb/\Ja = \dim_q(\tfrac{1}{2};k)$ satisfies
\begin{equation}
  \sopt^{2k} \;=\; N_{\mathrm{fus}}(Q,k),
  \label{P1:eq:lsc_p1}
\end{equation}
where $N_{\mathrm{fus}}(Q,k)$ is the number of fusion channels of
$\phi_{1/2}$ at level $k$.
Equivalently, $\leff^* = \lam/N_{\mathrm{fus}}^{1/k}$.
\end{conjecture}
\status{proved (Paper~II)}
\source{P1:conj:linking; proved as \texttt{conj:lsl} in Paper~II}


%════════════════════════════════════════════════════════════════════
% Saddle Point and BPS Bound — Paper II
%════════════════════════════════════════════════════════════════════
\begin{theorem}[Fixed-point value]
\label{P2:thm:main}
Subject to the conditions of Theorems~\ref{P2:thm:existence}
and~\ref{P2:thm:mono}, the self-consistency equation $H(\leff^*)=\leff^*$,
where $H(\leff):=\Kfb[f_{\leff}]/\Jfour[f_{\leff}]$, has a unique solution.
Under the Linking Scale Conjecture~\ref{P2:conj:lsl}, that solution satisfies
\begin{equation}
  \leff^* \;=\; \frac{\lam}{2^{1/3}} \;=\; \frac{\vp^6}{2^{1/3}}.
  \label{P2:eq:leffstar}
\end{equation}
Consequently, the EL minimum of $\Efb$ in the $Q=2$ sector satisfies
\begin{equation}
  \frac{\Jfour}{\Ja}\bigg|_{\mathrm{EL}}
  \;=\; \frac{2^{4/3}}{\vp^5}.
  \label{P2:eq:J4J2a}
\end{equation}
In the thin-torus approximation (restricting to the BS ansatz family),
equation~\eqref{P2:eq:J4J2a} is proved unconditionally
(Theorem~\ref{P2:thm:lsl_1d}).
For the full 2D EL equation, equation~\eqref{P2:eq:J4J2a} is proved
in Paper~III~\cite{Paper3} (Theorems~4.3 and~4.4 therein) via
three-fermion Witten bosonization and the Brouwer/Rellich--Kondrachov~\cite{Evans2010}
fixed-point argument.
\end{theorem}
\status{published}
\source{P2:thm:main}

\begin{theorem}[Existence, subject to Condition~\ref{P2:cond:ps}]
\label{P2:thm:existence}
Assume the following concentration-compactness condition holds:
\begin{condition}[Palais--Smale]
\label{P2:cond:ps}
Every sequence $\{f_n\}\subset X_2$ with $\G[f_n;\leff]\leq M$ and
$\|\delta\G[f_n;\leff]\|_{H^{-1}}\to 0$ has a strongly convergent
subsequence in $H^1_{\mathrm{loc}}$.
\end{condition}
Under Condition~\ref{P2:cond:ps}, for each $\leff>0$,
$\G[\,\cdot\,;\leff]$ attains a critical point $f_{\leff}\in X_2$.
\end{theorem}
\status{proved (subject to concentration-compactness; resolved in Paper~III)}
\source{P2:thm:existence}

\begin{theorem}[Monotonicity and uniqueness]
\label{P2:thm:mono}
The map $H\colon(0,\infty)\to(0,\infty)$ defined by
$H(\leff) = \Kfb[f_{\leff}]/\Jfour[f_{\leff}]$
satisfies:
\begin{enumerate}[label=\textup{(\roman*)}]
  \item $\dfrac{\mathrm{d}}{\mathrm{d}\leff}\bigl[H(\leff)-\leff\bigr] < 0$
    for all $\leff>0$, i.e., $H(\leff)-\leff$ is strictly decreasing.
  \item $H(\leff)-\leff > 0$ for $\leff \ll 1$ and
    $H(\leff)-\leff < 0$ for $\leff \gg 1$.
  \item The equation $H(\leff^*)=\leff^*$ has exactly one solution.
\end{enumerate}
\end{theorem}
\status{proved (subject to non-degeneracy of Hessian; existence fully resolved in Paper~III)}
\source{P2:thm:mono}

\begin{remark}[Non-degeneracy]
\label{P2:rem:nd}
The assumption that $\delta^2\G/\delta f^2$ is positive definite at
$f_{\leff}$ is equivalent to $f_{\leff}$ being a non-degenerate
local minimum of $\G$.
Numerically, gradient descent on $\G[\,\cdot\,;\leff]$ converges
stably for $\leff$ near $\leff^*$, with no zero modes observed.
Strict convexity of $\Kfb$ (since $x/(1+\bst x)$ is concave in $\kk$)
ensures $\delta^2\Kfb/\delta f^2 > 0$; together with the
Vakulenko--Kapitanski lower bound on $\Jfour$ in the $Q=2$ sector,
this gives $\delta^2\G/\delta f^2 > 0$ at the mountain-pass critical point.
\end{remark}
\status{proved (partial: Hessian spectral gap remains open; not required for main results)}
\source{P2:rem:nd}

\begin{theorem}[Linking Scale Lemma --- thin-torus case]
\label{P2:thm:lsl_1d}
In the thin-torus approximation, restrict to profiles of the form
$f(d) = 2\arctan(C/d)$ where $d = \sqrt{(r-R_0)^2+z^2}$ and $C>0$
is the tube-width parameter.  Define the 1D functionals
\begin{align}
  \Ia   &:= \int_0^\infty \sin^4\!f\cdot\kk\cdot 2\pi d\,\mathrm{d}d,
  &
  \Ib   &:= \int_0^\infty F_{13}^2\cdot 2\pi d\,\mathrm{d}d,
\end{align}
where $\kk = (f')^2 + \sin^2\!f/d^2$ and $F_{13} = -\sin^2\!f\cdot f'/d$
in the thin-torus limit $A\approx d^{-2}$.

\begin{enumerate}[label=\textup{(\roman*)}]
\item $\Ib/\Ia = 3/(4C^2)$ exactly.
\item The self-consistent fixed point satisfying $V=\vp$ and
  $\leff = \Kfb/\Ib$ has
  \[
    C^{*2} = \frac{3\vp^5}{8\cdot 2^{1/3}}, \qquad
    \frac{\Ib}{\Ia}\bigg|_{C=C^*} = \frac{2^{4/3}}{\vp^5}.
  \]
\end{enumerate}
\end{theorem}
\status{published}
\source{P2:thm:lsl_1d}

\begin{proposition}[S-matrix self-duality]
\label{P2:prop:Sselfdual}
$S_{\frac{1}{2},\frac{1}{2}}\big|_{k=3} = S_{0,0}\big|_{k=3}$.
\end{proposition}
\status{published}
\source{P2:prop:Sselfdual}

\begin{proposition}[Twist pairing]
\label{P2:prop:twist}
$\theta_0\cdot\theta_{\frac{1}{2}} = 1$ at $k=3$,
where $\theta_j = \exp\bigl(2\pi\mathrm{i}(h_j - c/24)\bigr)$.
\end{proposition}
\status{published}
\source{P2:prop:twist}

\begin{proposition}[Fusion count]
\label{P2:prop:Nfus}
$N_{\mathrm{fus}} := |\{\text{summands in }\tfrac{1}{2}{\times}\tfrac{1}{2}
  \text{ at level }k\}| = 2$ for all $k\geq 2$.
\end{proposition}
\status{published}
\source{P2:prop:Nfus}

\begin{proposition}[Sign identity]
\label{P2:prop:sign}
At a non-degenerate critical point $f_{\leff}$:
\begin{equation}
  \frac{\mathrm{d}H}{\mathrm{d}\leff}(\leff)
  \;=\;
  -\,\frac{\leff\Jfour + \Kfb}{\Jfour^2}
  \left\langle \frac{\delta\Jfour}{\delta f},
  \left(\frac{\delta^2\G}{\delta f^2}\right)^{-1}
  \frac{\delta\Jfour}{\delta f}\right\rangle.
  \label{P2:eq:sign}
\end{equation}
Since the inner product is strictly positive (positive-definiteness
of $\delta^2\G$), and $\leff\Jfour+\Kfb>0$, the right-hand side
is strictly negative: $\mathrm{d}H/\mathrm{d}\leff < 0$.
\end{proposition}
\status{published}
\source{P2:prop:sign}

\begin{conjecture}[Linking Scale --- proved in Paper~III~\cite{Paper3}]
\label{P2:conj:lsl}
For the density-feedback FN Hopfion at WZW level $k\geq 2$,
the unique fixed point of $H$ satisfies
\begin{equation}
  {\sopt}^{2k} \;=\; N_{\mathrm{fus}},
  \qquad
  N_{\mathrm{fus}}
  \;:=\;
  \bigl|\{\text{irreducible summands in }
  \tfrac{1}{2}\times\tfrac{1}{2}\text{ at level }k\}\bigr|,
  \label{P2:eq:lsl_conj}
\end{equation}
equivalently $\leff^* = \lam_k/N_{\mathrm{fus}}^{1/k}$.
\end{conjecture}
\status{proved (Paper~III)}
\source{P2:conj:lsl; proved in Paper~III}

\begin{equation}
  \frac{\delta\Kfb}{\delta f} = -\leff\,\frac{\delta\Jfour}{\delta f},
  \label{P2:eq:BPS}
\end{equation}

\begin{equation}
  \Jfour\,\frac{\delta\Kfb}{\delta f}
  + \Kfb\,\frac{\delta\Jfour}{\delta f} = 0.
  \label{P2:eq:EL}
\end{equation}

\begin{equation}
  \Efb[f] \;=\; \Kfb[f]\cdot\Jfour[f],
  \label{P2:eq:Efb}
\end{equation}

\begin{equation}
  \G[f;\leff] \;:=\; \Kfb[f] + \leff\,\Jfour[f].
  \label{P2:eq:G}
\end{equation}

\begin{equation}
  \frac{\mathrm{d}H}{\mathrm{d}\leff} < 0.
  \label{P2:eq:Hdecreasing}
\end{equation}

\begin{equation}
  \frac{\delta^2\G}{\delta f^2}\,\dot{f}
  = -\frac{\delta\Jfour}{\delta f}.
  \label{P2:eq:IFT}
\end{equation}

\begin{equation}
  \left(\frac{\lam}{\leff^*}\right)^k \;=\; N_{\mathrm{fus}}.
  \label{P2:eq:LSC_modular}
\end{equation}

\begin{equation}
  S_{\frac{1}{2},\frac{1}{2}}\big|_{k=3}
  \;=\; S_{0,0}\big|_{k=3},
  \label{P2:eq:Sselfdual}
\end{equation}

\begin{equation}
  \frac{\delta\Kfb}{\delta f} + \leff\,\frac{\delta\Jfour}{\delta f} = 0.
  \label{P2:eq:crit}
\end{equation}

\begin{equation}
  \frac{\mathrm{d}H}{\mathrm{d}\leff}
  = \frac{1}{\Jfour^2}
    \left[
      \left\langle \frac{\delta\Kfb}{\delta f}, \dot{f}\right\rangle \Jfour
      - \Kfb\left\langle \frac{\delta\Jfour}{\delta f}, \dot{f}\right\rangle
    \right].
  \label{P2:eq:dHdlam}
\end{equation}

\begin{align}
  \frac{\mathrm{d}H}{\mathrm{d}\leff}
  &= \frac{1}{\Jfour^2}
    \left[
      -\leff\left\langle \frac{\delta\Jfour}{\delta f},\dot{f}\right\rangle\Jfour
      - \Kfb\left\langle \frac{\delta\Jfour}{\delta f},\dot{f}\right\rangle
    \right] \notag\\
  &= -\frac{\leff\Jfour + \Kfb}{\Jfour^2}
    \left\langle \frac{\delta\Jfour}{\delta f}, \dot{f}\right\rangle.
  \label{P2:eq:dHdlam2}
\end{align}

\begin{equation}
  \frac{\Jfour}{\Ja}\bigg|_{\mathrm{EL}}
  \;=\; \frac{N_{\mathrm{fus}}^{1+1/k}}{[\dim_q(\tfrac{1}{2};k)]^{2k-1}}.
  \label{P2:eq:general_J4J2a}
\end{equation}

\begin{equation}
  \leff^* \;=\; \frac{\Kfb}{\Jfour} \;=\; \frac{2\vp\Ja}{\Jfour}
  \;=\; \frac{2\vp}{\Jfour/\Ja}.
  \label{P2:eq:leff_from_ratio}
\end{equation}

\begin{align}
  \leff^* &= \frac{\lam}{2^{1/k}}\bigg|_{k=3} = \frac{\lam}{2^{1/3}},
  \label{P2:eq:leffstar_v1}\\
  \frac{\Jfour}{\Ja}\bigg|_{\mathrm{EL}} &= \frac{2^{4/3}}{\vp^5},
  \label{P2:eq:leffstar_v2}\\
  {\sopt}^2 &= 2^{1/3} \qquad\text{(i.e.\ }\sopt = 2^{1/6}\text{)},
  \label{P2:eq:leffstar_v3}\\
  {\sopt}^{2k} &= 2 \qquad\text{at }k=3.
  \label{P2:eq:leffstar_v4}
\end{align}

\begin{equation}
  \left\langle \frac{\delta\Jfour}{\delta f}, \dot{f}\right\rangle
  = -\left\langle \frac{\delta\Jfour}{\delta f},
    \left(\frac{\delta^2\G}{\delta f^2}\right)^{-1}
    \frac{\delta\Jfour}{\delta f}\right\rangle.
  \label{P2:eq:quad}
\end{equation}


%════════════════════════════════════════════════════════════════════
% Fixed Point and Virial Algebra — Paper III (selected)
%════════════════════════════════════════════════════════════════════
\begin{proposition}[Exact thin-torus virial function]
\label{P3:prop:virial_exact}
For the BS profile $f = 2\arctan(C/\rho)$ in the thin-torus limit,
the virial ratio $V(\beta) = \Jfb/\Ja$ has the closed-form expression
\begin{equation}
  V(b) \;=\; \frac{15}{8} \cdot \frac{\arctan\!\sqrt{8b}}{\sqrt{8b}},
  \qquad b = \beta/C^2,
  \label{P3:eq:V_exact}
\end{equation}
for all $C > 0$ and $\beta \geq 0$.
\end{proposition}
\status{published}
\source{P3:prop:virial_exact}

\begin{lemma}[Virial algebra]
\label{P3:lem:virial_alg}
$1 + \mu^*\vp = 2\vp$.
\end{lemma}
\status{published}
\source{P3:lem:virial_alg}

\begin{lemma}[Golden identity]
\label{P1:lem:golden}
$1 + \ms\vp = 2\vp$ exactly.
\end{lemma}
\status{published}
\source{P1:lem:golden}

\begin{corollary}[Existence of the feedback fixed point]
\label{P3:cor:fixedpoint_exists}
For any $R_0$, there exists a unique $b^*(R_0) > 0$ such that
$V(b^*) = \vp$.  Numerically, $b^* \approx 0.06715$ and the self-consistent
tube width satisfies $C^{*2}(R_0) = 3/[4(\mathrm{WZW} - 1/(2R_0^2))]$.
\end{corollary}
\status{published}
\source{P3:cor:fixedpoint_exists}

\begin{remark}[Structure of the fixed-point equation]
\label{P3:rem:fixedpoint_structure}
The condition $V(b^*) = \vp$ is the WZW self-consistency equation
for the density-feedback model.  In renormalization-group language:
the coupling ratio $V$ runs from the ultraviolet value $V(0) = 15/8$
to the infrared fixed point $V(b^*) = \vp = \dim_q(\tfrac{1}{2};k{=}3)$
under the arctan flow~\eqref{P3:eq:V_exact}.  The UV starting value
$15/8 = I_{2\mathrm{iso}}/I_{2a} = 4/(32/15)$ is determined by the ratio
of two C-independent integrals of the BS profile; both C-independence
facts follow from the key identity $\sin^2\!f/\rho^2 = (f')^2$
(Proposition~\ref{P3:prop:scale_inv}).
\end{remark}
\status{published}
\source{P3:rem:fixedpoint_structure}


%════════════════════════════════════════════════════════════════════
% LSC Proof Chain — Paper III
%════════════════════════════════════════════════════════════════════

\begin{conjecture}[Linking Scale Conjecture --- proved in Paper~III]
\label{P3:conj:lsc}
The Euler--Lagrange minimum of $\Efb = \Kfb\cdot\Jfour$ in the $Q=2$
topological sector satisfies
\begin{equation}
  \frac{\Jfour}{\Ja}\bigg|_{\mathrm{EL}}
  \;=\; \frac{2^{4/3}}{\vp^5}
  \;\approx\; 0.22721,
  \label{P3:eq:lsc}
\end{equation}
equivalently $\leff^* = \lam/2^{1/3}$ where $\lam = \vp^6$.
\end{conjecture}
\status{proved (Paper~III via Theorem~\ref{P3:thm:V_exact_all_R0} + Theorem~\ref{P3:thm:lsc_conditional})}
\source{P3:conj:lsc}

\begin{lemma}[Bridge identity]
\label{P3:lem:bridge}
The Morse parameter $\mu^* = 3-\vp$ of the density-feedback
functional satisfies
\begin{equation}
  \mu^* \;=\; \frac{\sqrt{k+2}}{\vp}
  \label{P3:eq:bridge}
\end{equation}
with $k=3$.
\end{lemma}
\status{published}
\source{P3:lem:bridge}
\residual{exact}

\begin{remark}
\label{P3:rem:bridge_interpretation}
Lemma~\ref{P3:lem:bridge} is the algebraic bridge connecting the
condensate parameter $\mu^*$ to the WZW level $k=3$.
The coefficient $(k+2)=5$ appears in the Knizhnik--Zamolodchikov
equation~\cite{KZ1984} for $\mathrm{SU}(2)_k$ and here emerges
from the icosahedral geometry of $\mu^*$.
\end{remark}
\status{published}
\source{P3:rem:bridge_interpretation}

\begin{theorem}[LSC conditional on virial]
\label{P3:thm:lsc_conditional}
Suppose the Euler--Lagrange minimum of $E_{\mathrm{fb}}$ satisfies
the virial condition
\begin{equation}
  V^* \;:=\; \frac{\Jfb}{\Ja}\bigg|_{\mathrm{EL}} \;=\; \vp.
  \label{P3:eq:virial}
\end{equation}
Then $\Jfour/\Ja = 2^{4/3}/\vp^5$, i.e.\ the LSC holds.
\end{theorem}
\status{published}
\source{P3:thm:lsc_conditional}
\begin{proof}
$\Kfb = \Ja(1 + \mu^*V^*) = \Ja(1 + \mu^*\vp) = 2\vp\Ja$
(using Lemma~\ref{P3:lem:virial_alg}).
Then $\Jfour/\Ja = \Kfb/(\Ja\leff^*) = 2\vp/(\vp^6/2^{1/3}) = 2^{4/3}/\vp^5$.
\end{proof}
\status{published}
\source{P3:thm:lsc_conditional}
\residual{exact}

\begin{remark}
\label{P3:rem:lsc_conditional_status}
Theorem~\ref{P3:thm:lsc_conditional} reduces the LSC to the virial
condition $V^*=\vp$.  Section~\ref{P3:sec:bosonization} proves $V^*=\vp$
exactly (thin-torus via Verlinde; all $R_0$ via Theorem~\ref{P3:thm:V_exact_all_R0}),
so the LSC is a theorem.
\end{remark}
\status{published}
\source{P3:rem:lsc_conditional_status}

\begin{theorem}[Equivalence of LSC, virial, and fixed-point conditions]
\label{P3:thm:lsc_equivalence}
At the EL minimum of $\Efb$ in the $Q=2$ sector, the following are
pairwise equivalent:
\begin{enumerate}
  \item[(A)] $V^* = \Jfb/\Ja = \vp$
  \item[(B)] $\Jfour/\Ja = 2^{4/3}/\vp^5$
  \item[(C)] $\leff^* = \lam/2^{1/3}$
\end{enumerate}
\end{theorem}
\status{published}
\source{P3:thm:lsc_equivalence}
\residual{exact}

\begin{remark}[All three conditions are proved]
\label{P3:rem:all_three_proved}
Theorem~\ref{P3:thm:lsc_equivalence} shows that conditions (A), (B), (C) are
the same statement.  All three are proved: $V^*=\vp$ by the
three-fermion bosonization (Section~\ref{P3:sec:bosonization},
Theorem~\ref{P3:thm:V_exact_all_R0}), from which (B) and (C) follow
immediately by the equivalences established above.
\end{remark}
\status{published}
\source{P3:rem:all_three_proved}

\begin{proposition}[Scale invariance of thin-torus functionals]
\label{P3:prop:scale_inv}
In the thin-torus approximation with profile $f = 2\arctan(C/\rho)$,
\begin{align}
  J_{2a} &= 2\pi R_0 \cdot \frac{16}{15}, \label{P3:eq:Ja_exact}\\
  J_4    &= 2\pi R_0 \cdot \frac{64^2}{2}\,B(5,7), \label{P3:eq:J4_exact}
\end{align}
where $B(5,7) = \Gamma(5)\Gamma(7)/\Gamma(12)$.  Both are independent of $C$.
\end{proposition}
\status{published}
\source{P3:prop:scale_inv}
\residual{exact}

\begin{remark}
\label{P3:rem:scale_invariance_explains}
Scale invariance explains why $J_4/J_{2a} = 2^{4/3}/\vp^5$ is a pure
number independent of $C$.  The self-consistent $C^*$ is fixed by the
EL equation, not by the ratio $J_4/J_{2a}$.
\end{remark}
\status{published}
\source{P3:rem:scale_invariance_explains}

\begin{proposition}[Thin-torus virial limit]
\label{P3:prop:virial_limit}
$V(0;R_0) \to \vp$ as $R_0\to\infty$, with the leading correction
\begin{equation}
  V(0;R_0) \;=\; \vp \;+\; \frac{A}{R_0^2} \;+\; O(R_0^{-4}),
  \label{P3:eq:V0_limit}
\end{equation}
for a computable coefficient $A > 0$.
\end{proposition}
\status{published}
\source{P3:prop:virial_limit}

\begin{remark}[Exact for all $R_0$]
\label{P3:rem:exact_for_all_R0}
Proposition~\ref{P3:prop:virial_limit} shows that $\vp$ is the exact
continuum-limit attractor for $V(0;R_0)$, approached as $1/R_0^2$.
The thin-torus self-consistent profile satisfies $V^*=\vp$ exactly
(Theorems~\ref{P3:thm:lsc_1d} and~\ref{P3:thm:lsc_conditional}).
Theorem~\ref{P3:thm:V_exact_all_R0} extends this to all finite $R_0$
via a fixed-point existence argument: $V^*=\vp$ holds exactly
throughout, not merely asymptotically.
\end{remark}
\status{published}
\source{P3:rem:exact_for_all_R0}

\begin{proposition}[Hypothesis (i) at large $R_0$]
\label{P3:prop:V0_gt_phi}
From Proposition~\ref{P3:prop:virial_limit}:
$V_0(R_0) = \vp + A/R_0^2 + O(R_0^{-4})$ with $A > 0$.
Hence $V_0(R_0) > \vp$ for all $R_0 < \infty$.
\end{proposition}
\status{published}
\source{P3:prop:V0_gt_phi}

\begin{theorem}[Universal virial identity --- thin-torus]
\label{P3:conj:kappa_universal}
For the $Q=2$ Hopfion in the self-consistent feedback regime, the
kappa-distribution $P_{R_0}(\kk)$ satisfies
$\int_0^\infty P_{R_0}(\kk)/(1+\bst(R_0)\kk)\,d\kk = \vp$
in the thin-torus limit.  The value $\vp = \dim_q(\tfrac{1}{2};k{=}3)$
is proved via three-fermion bosonization.
\end{theorem}
\status{proved (via Theorem~\ref{P3:thm:bos_main})}
\source{P3:conj:kappa_universal}

\begin{theorem}[Virial fixed-point theorem]
\label{P3:thm:virial_fixedpt}
Suppose (i) $V_0(R_0) > \vp$, (ii) $V(f^*(\beta),\beta)$ is strictly
decreasing and continuous in $\beta\geq 0$, (iii) $V\to 0$ as $\beta\to\infty$.
Then there exists a unique $\beta^*(R_0)>0$ such that $V(f^*(\beta^*),\beta^*)=\vp$,
and Theorem~\ref{P3:thm:lsc_conditional} gives $\Jfour/\Ja = 2^{4/3}/\vp^5$.
\end{theorem}
\status{proved (hypotheses (i) by Prop.~\ref{P3:prop:V0_gt_phi}; (ii) numerically verified; (iii) trivial)}
\source{P3:thm:virial_fixedpt}

\begin{remark}[Status of the hypotheses]
\label{P3:rem:hypotheses_status}
Hypothesis~(iii) holds trivially: for large $\beta$, the feedback suppresses all
curvature and $V\to 0$.
Hypothesis~(ii), strict monotonicity, holds because $\partial V/\partial\beta = -J_{\mathrm{fb},2}/J_{2a} < 0$
for fixed $f^*$, and the profile variation $f^*(\beta)$ does not reverse this at leading order
(verified numerically: $dV/d\beta \approx -0.82/J_{2a}$ at the fixed point).
Hypothesis~(i) is equivalent to: $J_{2\mathrm{iso}}(f^*_0) > \vp \cdot J_{2a}(f^*_0)$.
\end{remark}
\status{published}
\source{P3:rem:hypotheses_status}

\begin{remark}[Relation to Theorem~\ref{P3:thm:V_exact_all_R0}]
\label{P3:rem:relation_to_V_exact}
Theorem~\ref{P3:thm:virial_fixedpt} above gives $V^*=\vp$ from three
hypotheses; hypothesis~(ii) (strict monotonicity) is verified numerically
but not proved analytically by this route.
The stronger result --- $V^*=\vp$ exactly for all finite $R_0$ without
requiring hypothesis~(ii) --- is proved directly in
Theorem~\ref{P3:thm:V_exact_all_R0} via Brouwer's fixed-point theorem and
Rellich--Kondrachov compactness~\cite{Evans2010}.
Theorem~\ref{P3:thm:virial_fixedpt} is retained here because its three-hypothesis
structure gives a transparent account of the feedback mechanism.
\end{remark}
\status{published}
\source{P3:rem:relation_to_V_exact}

\begin{remark}[Curvature decomposition and the WZW metric]
\label{P3:rem:kappa_wzw_preview}
The isotropic curvature splits as $\kk_{\mathrm{iso}} = \kk_{\mathrm{tube}} + \sin^2\!f/r^2$.
The BS identity $\sin^2\!f/\rho^2 = (f')^2$ gives $\kk_{\mathrm{tube}} = 2(f')^2 = g_{\mathrm{WZW}}$.
At $R_0=3$, $\bst=0.452$: $V^* = 1.61824$ ($\vp=1.61803$, gap $0.013\%$) and
$\Jfour/\Ja = 0.22717$ ($2^{4/3}/\vp^5=0.22721$, gap $0.020\%$), both lattice artefacts.
\end{remark}
\status{published}
\source{P3:rem:kappa_wzw_preview}

\begin{theorem}[Thin-torus LSC --- \cite{Paper2}]
\label{P3:thm:lsc_1d}
In the thin-torus approximation with profiles $f = 2\arctan(C/d)$,
one has $I_4/I_{2a} = 3/(4C^2)$ exactly, from which self-consistency
gives $I_4/I_{2a}\big|_{C=C^*} = 2^{4/3}/\vp^5$.
\end{theorem}
\status{published}
\source{P3:thm:lsc_1d}

\begin{theorem}[Thin-torus LSC: complete proof]
\label{P3:thm:lsc_1d_complete}
For $f = 2\arctan(C/\rho)$: $I_{2a} = 32/15$, $I_{4\text{-tube}} = 8/(5C^2)$,
$I_{4\text{-wind}} = 16/15$ (all exact, $C$-independent via BS identity).
The full ratio is $J_4/J_{2a} = 3/(4C^2) + 1/(2R_0^2)$.
Given $V^*=\vp$, the WZW master relation forces $C^{*2}$ via
$3/(4C^{*2}) = 2^{4/3}/\vp^5 - 1/(2R_0^2)$,
after which the winding cancellation gives $J_4/J_{2a} = 2^{4/3}/\vp^5$ exactly.
\end{theorem}
\status{published}
\source{P3:thm:lsc_1d_complete}
\residual{exact}

\begin{remark}[Structural content of the cancellation]
\label{P3:rem:cancellation_structure}
The exact cancellation~\eqref{P3:eq:cancellation} rests on two
independent C-invariances: $I_{2a} = 32/15$ and $I_{4\text{-wind}} = 16/15$.
Both follow from the BS identity $\sin^2\!f/\rho^2 = (f')^2$,
which in turn follows from the specific profile form $f=2\arctan(C/\rho)$.
The ratio $I_{4\text{-wind}}/I_{2a}=1/2$ is thus a consequence of
the self-similar structure of the BS profile, not of $\vp$ or the WZW level.
The WZW level enters through the master relation~\eqref{P3:eq:master_relation},
which fixes $3/(4C^{*2}) = \text{WZW} - 1/(2R_0^2)$ via $V^*=\vp$.
\end{remark}
\status{published}
\source{P3:rem:cancellation_structure}

\begin{theorem}[LSC in the continuum limit]
\label{P3:thm:lsc_2d_limit}
If $V^*(R_0) = \vp + O(R_0^{-2})$ as $R_0\to\infty$, then
$J_4/J_{2a}\big|_{f^*_{R_0}} = 2^{4/3}/\vp^5 + O(R_0^{-2})$.
In particular $J_4/J_{2a} \to 2^{4/3}/\vp^5$ as $R_0\to\infty$.
\end{theorem}
\status{published}
\source{P3:thm:lsc_2d_limit}

\begin{theorem}[Exact $V^*=\vp$ for all finite $R_0$]
\label{P3:thm:V_exact_all_R0}
For each $R_0>0$, the density-feedback model has a self-consistent state
$(f^*,\bst)$ satisfying the EL equation of $\Kfb\cdot\Jfour$ with
$V(f^*,\bst) = \vp$ exactly.
\end{theorem}
\status{published}
\source{P3:thm:V_exact_all_R0}
\begin{proof}[Proof sketch]
The map $T: f\mapsto f^*(\beta^*(f))$ (EL minimiser at self-consistent $\beta$)
acts on bounded closed convex $Q=2$ profiles.
On finite grids, Brouwer's fixed-point theorem gives a fixed point $f^*_h$.
The sequence $\{f^*_h\}$ is $H^1$-bounded; Rellich--Kondrachov gives a
convergent subsequence whose limit is a continuum fixed point with $V^*=\vp$.
\end{proof}
\status{published}
\source{P3:thm:V_exact_all_R0}
\residual{exact (lattice artefact 0.013\% at $h=0.12$)}
\depends{profile/virial.tex (P3:thm:lsc_conditional, P3:lem:virial_alg)}

\begin{remark}[The 2D/thin-torus distinction, resolved]
\label{P3:rem:2d_thintorus_resolved}
Theorem~\ref{P3:thm:V_exact_all_R0} shows $V^* = \vp$ holds exactly at finite $R_0$,
not merely asymptotically.  The $O(R_0^{-2})$ statement in
Proposition~\ref{P3:prop:virial_limit} describes how $V(f,0)$ approaches $\vp$
as $R_0\to\infty$ with \emph{no feedback} ($\beta=0$); the self-consistent
$\bst(R_0) > 0$ then adjusts to drive $V$ exactly to $\vp$ for every $R_0$.
The numerical residual of $0.013\%$ at $R_0=3$ is a finite-grid artefact
($h=0.12$), not a theoretical gap.  The residual open problem is the
Hessian spectral gap condition ensuring the fixed point is unique
(Section~\ref{P3:sec:open}), which is not required for the LSC or $\alpha$ results.
\end{remark}
\status{published}
\source{P3:rem:2d_thintorus_resolved}

\begin{remark}[Role of $\mu^* = \sqrt{k+2}/\vp$]
\label{P3:rem:mustar_role}
Lemma~\ref{P3:lem:bridge} now acquires a new interpretation.
$\mu^*=\sqrt{k+2}/\vp$, evaluated at $k=3$ --- the Wess--Zumino coupling
constant (level) of the $\mathrm{SU}(2)_3$ theory in the standard sense
of~\cite{Witten1984} --- gives the value $3-\vp$ that appears in the
density-feedback self-consistency condition.
$\mu^*$ itself is \emph{not} a coupling constant in the WZW sense (that
role is played by $k$, which is quantised to an integer by the standard
topological argument); $\mu^*$ is the feedback-model parameter whose
value happens to be pinned down once $k$ is fixed.
The density-feedback model with $\mu = \mu^*$ is the \emph{unique} value
at which $K_{\mathrm{fb}}$ is the WZW action; any other $\mu$ gives a
deformation away from the fixed point.  At $\mu = \mu^*$ the WZW Verlinde
formula applies and delivers $V^* = \vp$ exactly.
\end{remark}
\status{published}
\source{P3:rem:mustar_role}

\begin{proposition}[Resolvent-to-primary identification]
\label{P3:prop:res_primary}
In the thin-torus limit, $\kk_{\mathrm{tube}} = 2(f')^2 = g_{\mathrm{WZW}}$
(the WZW metric), and the feedback density $\kk/(1+\beta\kk)$ is the
Stieltjes resolvent of the WZW spectral measure.  At $\beta=\bst$:
$V(\bst) = (15/8)\arctan(\sqrt{8\bst})/\sqrt{8\bst} = \vp = S_{0,1/2}/S_{0,0}$.
\end{proposition}
\status{published}
\source{P3:prop:res_primary}
\residual{exact}

\begin{remark}[2D corrections]
\label{P3:rem:2d_corrections}
For finite $R_0$, the full curvature $\kappa_{\mathrm{iso}} = \kappa_{\mathrm{tube}}
+ \sin^2\!f/r^2$ includes a winding term $\sin^2\!f/r^2$ absent from the WZW metric.
This deforms $V$ by $O(R_0^{-2})$, shifting $\bst(R_0)$ accordingly.
The self-consistent virial satisfies $V^*(R_0) = \vp + O(R_0^{-2})$
(Proposition~\ref{P3:prop:virial_limit}).  The exact cancellation in $J_4/J_{2a}$
(Theorem~\ref{P3:thm:lsc_1d_complete}) shows the winding correction fills the tube
deficit precisely, giving $J_4/J_{2a} = 2^{4/3}/\vp^5$ exactly for all $R_0$.
\end{remark}
\status{published}
\source{P3:rem:2d_corrections}

%════════════════════════════════════════════════════════════════════
% Universal Virial Identity (all sectors) — Paper XV
%════════════════════════════════════════════════════════════════════

\begin{theorem}[Universal virial identity]
\label{P15:thm:virial}
At the density-feedback virial saddle of any sector $n$,
$V^{*\,(n)}\equiv J_{\mathrm{fb}}^{(n)}/\Ja^{(n)}=\vp$,
and consequently
\begin{equation}
  \Kfb^{(n)} = 2\vp\,\Ja^{(n)},
  \qquad
  \Efb^{(n)} = 2N_n^2/\vp^{17}.
  \label{P15:eq:virial_universal}
\end{equation}
The integer $N_n = Q_{\mathrm{group}}^{(n)}$ is the T-matrix closing
integer of the sector. The formula is proved via the IVT argument
of Paper~III applied sector-independently to the trefoil geometry.
\end{theorem}
\status{published}
\source{P15:thm:virial}
\residual{exact}

\begin{remark}[Sector independence of the energy formula]
\label{P15:rem:sector_independence_energy}
Equation~\eqref{P15:eq:virial} shows that $\Efb^{(n)}=2N_n^2/\vp^{17}$
depends only on the integer $N_n$, not on the virial ratio $r_n$.
The sector geometry changes $r_n$ but not the overall energy scale
(at fixed $N_n$).
\end{remark}
\status{published}
\source{P15:rem:sector_independence_energy}
